% =============================================================================
% Partie 2 — Problématique et objectifs
% =============================================================================
\chapter{Problématique et objectifs}

\epigraphe{Un problème bien posé est à moitié résolu.}

Ta responsable t'a donné quatre problèmes à résoudre. Avant de les traduire en SQL, il faut les
comprendre parfaitement : ce que mesure chaque indicateur, pourquoi la définition exacte compte, et ce
que la banque risque si on se trompe. Cette partie les pose un par un, avec des calculs à la main sur
de vrais prêts de ta base.

% -----------------------------------------------------------------------------
\section{Les quatre problèmes métier à résoudre}

\subsection{Problème 1 --- Suivre les retards de paiement (DPD)}

\begin{imagine}
Deux clients ont chacun \textbf{une} mensualité impayée. Le premier a oublié celle d'hier. Le second
n'a pas payé celle d'il y a quatre mois, mais il a payé normalement les trois suivantes.

Sont-ils dans la même situation ? Évidemment pas : le premier a un jour de retard, le second plus de
cent. Compter les impayés ne suffit pas : il faut savoir \textbf{depuis quand} le client doit de
l'argent.
\end{imagine}

\begin{definition}{Days Past Due (DPD), le nombre de jours de retard}
À une date d'arrêté $t$, le DPD d'un prêt est le nombre de jours écoulés depuis \textbf{l'échéance la
plus ancienne qui est échue et pas entièrement payée}. S'il n'y en a aucune, le DPD vaut 0.
\[
DPD(t) = t - \min\{\, d_k \;:\; d_k \leq t \text{ et l'échéance } k \text{ n'est pas soldée à } t \,\}
\]
(et $DPD(t) = 0$ si aucune échéance ne remplit ces conditions),
où $d_k$ est la date de l'échéance numéro $k$.
\end{definition}

\begin{calcul}{le retard du prêt 61 au 30/09/2026}
Voici les mensualités du prêt 61 autour de la date d'arrêté (704,41~€ par mois) :

\begin{center}\small
\begin{tabular}{@{}cccccl@{}}
\toprule
N° & Échéance & Montant dû & Payé & Date de paiement & État au 30/09 \\
\midrule
15 & 15/05/2026 & 704,41 & 704,41 & 18/05/2026 & soldée \\
16 & 15/06/2026 & 704,41 & 704,41 & 15/06/2026 & soldée \\
17 & 15/07/2026 & 704,41 & 704,41 & 03/08/2026 & soldée (avec 19 jours de retard) \\
\rowcolor{rougeFond} 18 & 15/08/2026 & 704,41 & 0,00 & --- & \textbf{échue, impayée} \\
19 & 15/09/2026 & 704,41 & 704,41 & 15/09/2026 & soldée \\
20 & 15/10/2026 & 704,41 & 0,00 & --- & pas encore échue \\
\bottomrule
\end{tabular}
\end{center}

\begin{enumerate}
\item Échéances échues au 30/09 : les numéros 15 à 19 (la 20\textsuperscript{e} est dans le futur).
\item Parmi elles, non soldées : seulement la \textbf{18} (la 19 a été payée le 15/09).
\item La plus ancienne non soldée est donc due depuis le \textbf{15/08/2026}.
\item $DPD = 30/09/2026 - 15/08/2026 = 16 + 30 = \mathbf{46}$ \textbf{jours}.
\end{enumerate}
Le prêt 61 a 46 jours de retard, bien qu'il ait payé sa dernière mensualité à l'heure.
\end{calcul}

\begin{piege}{Piège --- Trois fausses définitions du retard}
\begin{itemize}
\item \textbf{« Jours depuis le dernier paiement »} : pour le prêt 61, cela donnerait 15 jours (dernier
paiement le 15/09). Faux : l'impayé d'août est masqué par le paiement de septembre.
\item \textbf{« Nombre d'impayés × 30 »} : une mensualité impayée donnerait 30 jours. Faux : les mois
n'ont pas tous 30 jours et l'impayé date de 46 jours.
\item \textbf{« Jours depuis l'impayé le plus récent »} : avec plusieurs impayés, on sous-estimerait
gravement le retard. La règle bancaire retient toujours le \textbf{plus ancien}.
\end{itemize}
\end{piege}

\begin{observation}[Pourquoi le DPD n'est pas une colonne de la base]
Le retard change \textbf{tous les jours}, sans qu'aucun événement ne se produise : un prêt qui a 46
jours de retard aujourd'hui en aura 47 demain. Une colonne \code{jours_retard} stockée serait donc
fausse dès le lendemain, sauf à la recalculer chaque nuit pour tous les prêts. On la
\textbf{calcule} à partir des dates d'échéance et des paiements, dans une vue. C'est un principe général
de conception : \emph{une donnée calculée et stockée finit toujours par devenir incohérente}.
\end{observation}

Pour piloter un portefeuille, on regroupe les prêts par \textbf{tranche de retard} (\emph{buckets}) :
à jour, 1 à 30 jours, 31 à 60, 61 à 90, plus de 90 jours. On en déduit deux indicateurs suivis par
tous les comités des risques.

\begin{definition}{PAR30 et PAR90 (Portfolio At Risk)}
Le PAR$x$ est la part de l'encours (le capital restant dû) portée par les prêts qui ont \textbf{plus
de $x$ jours} de retard :
\[
PAR_x = \frac{\text{encours des prêts dont } DPD > x}{\text{encours total}}
\]
On le calcule en \textbf{montant}, pas en nombre de prêts : un gros prêt en retard pèse plus qu'un
petit.
\end{definition}

\subsection{Problème 2 --- Calculer l'exposition (EAD) et les provisions}

Pour un prêt amortissable, l'\textbf{exposition} (EAD) est le capital que le client doit encore. Il
diminue à chaque mensualité : une mensualité contient une part de capital et une part d'intérêts, et
seule la part de capital réduit la dette. Si une mensualité n'est payée qu'en partie, on considère que
le paiement couvre capital et intérêts au prorata.

La provision IFRS~9 d'un prêt est sa perte attendue, avec une probabilité de défaut dont
\textbf{l'horizon dépend du stage} :

\begin{center}\small
\begin{tabular}{@{}cllc@{}}
\toprule
Stage & Condition (dans ce projet) & Horizon de la PD & PD utilisée \\
\midrule
1 & $DPD \leq 30$ & 12 prochains mois & $PD_{12}$ \\
2 & $30 < DPD \leq 90$ & toute la durée restante & $PD_{\text{restante}}$ \\
3 & $DPD > 90$ (défaut) & --- & 100~\% \\
\bottomrule
\end{tabular}
\end{center}
\[
\boxed{\;ECL = PD_{\text{stage}} \times LGD \times EAD\;}
\]

\subsubsection{De la PD du modèle à la PD sur 12 mois : la dérivation}

Ton modèle prédit une PD sur \textbf{toute la durée} du prêt (la variable cible de ton jeu de données
indique si le prêt a fait défaut pendant sa vie). Or le stage 1 demande une PD sur 12 mois. Comment
passer de l'une à l'autre ?

\begin{analogie}{la pièce truquée}
Imagine qu'à chaque fin de mois, on lance une pièce truquée qui tombe sur « défaut » avec une petite
probabilité $h$, toujours la même. Le prêt survit s'il évite « défaut » à chaque lancer. Plus le prêt
dure, plus il y a de lancers, plus le risque cumulé augmente.
\end{analogie}

On fait donc l'hypothèse d'un \textbf{risque mensuel constant} $h$. Pour un prêt de $n$ mois :
\begin{enumerate}
\item Le prêt ne fait jamais défaut s'il survit aux $n$ mois : $P(\text{survie}) = (1-h)^n$.
\item Cette probabilité vaut $1 - PD$, où $PD$ est la PD du modèle sur la durée totale :
\[ (1-h)^n = 1 - PD \quad\Longrightarrow\quad 1 - h = (1 - PD)^{1/n} \]
\item La probabilité de faire défaut au cours des $m$ prochains mois est donc :
\[ PD_m = 1 - (1-h)^m = 1 - \left((1 - PD)^{1/n}\right)^m \]
\end{enumerate}
\[
\boxed{\;PD_m = 1 - (1 - PD)^{m/n}\;}
\]
Pour le stage 1, $m = \min(\text{mois restants}, 12)$ ; pour le stage 2, $m$ = mois restants.

\begin{observation}[Une hypothèse simplificatrice, à assumer en entretien]
Les banques utilisent des « structures par terme » de PD estimées sur leur historique : le risque n'est
pas constant (il est souvent plus élevé dans les premiers mois). L'hypothèse du risque constant est
simple, transparente et suffisante pour ce projet. Savoir dire \emph{pourquoi} c'est une approximation
vaut mieux qu'un modèle compliqué mal compris.
\end{observation}

\begin{calcul}{les provisions de trois prêts au 30/09/2026}
LGD = 45~\% pour tous.
\begin{itemize}
\item \textbf{Prêt 67} --- 3 jours de retard : \textbf{stage 1}. EAD = 16\,680,72~€, PD du modèle
0,0973~\%, durée 36 mois, 19 mois restants, donc $m = \min(19, 12) = 12$.
\[ PD_{12} = 1 - (1 - 0{,}000973)^{12/36} = 0{,}0324~\% \]
\[ ECL = 0{,}000324 \times 0{,}45 \times 16\,680{,}72 = \mathbf{2{,}43~€} \]
\item \textbf{Prêt 61} --- 46 jours de retard : \textbf{stage 2}. EAD = 19\,578,01~€, PD 1,0725~\%,
durée 48 mois, 29 mois restants.
\[ PD_{29} = 1 - (1 - 0{,}010725)^{29/48} = 0{,}6494~\% \]
\[ ECL = 0{,}006494 \times 0{,}45 \times 19\,578{,}01 = \mathbf{57{,}21~€} \]
\item \textbf{Prêt 9} --- 388 jours de retard : \textbf{stage 3}, $PD = 100~\%$. EAD = 21\,712,35~€.
\[ ECL = 1 \times 0{,}45 \times 21\,712{,}35 = \mathbf{9\,770{,}56~€} \]
\end{itemize}
Le stage pèse bien plus que la note : le prêt 9 avait une bonne note A, mais il est en défaut.
\end{calcul}

\subsection{Problème 3 --- Historiser les prédictions des modèles}

Le scoring ne se fait pas une fois pour toutes. Une banque fait tourner plusieurs modèles, les
réentraîne, change leurs seuils. Pour chaque prêt, il faut pouvoir répondre à des questions comme :
\emph{quel modèle a décidé ? avec quel seuil ? quand ? pour quelle raison ? les autres modèles
auraient-ils décidé autrement ?}

Voici, pour le prêt 61, les trois prédictions enregistrées dans ta base :

\begin{center}\small
\begin{tabular}{@{}lcrrcl@{}}
\toprule
Modèle & Champion & PD & Seuil & Décision & Explication principale (SHAP) \\
\midrule
\code{grille_score_v1} & non & 3,69~\% & 24,0~\% & accorder & --- (623 points) \\
\code{logistique_v1} & non & 2,66~\% & 23,5~\% & accorder & --- \\
\code{xgboost_monotone_v2} & \textbf{oui} & 1,07~\% & 25,5~\% & accorder & \code{note_risque} : $-1{,}29$ \\
\bottomrule
\end{tabular}
\end{center}

Cette table est un \textbf{registre} : chaque ligne est une prédiction d'un modèle pour un prêt, et elle
ne sera jamais modifiée. Ce registre sert à quatre choses :
\begin{enumerate}
\item \textbf{L'audit} : retrouver la version, le seuil et la date qui ont conduit à la décision.
\item \textbf{L'explication} : répondre à un client refusé, grâce à la variable SHAP principale.
\item \textbf{Le monitoring} : comparer, des mois plus tard, la PD prédite au défaut observé, et
détecter qu'un modèle se dégrade.
\item \textbf{Le champion contre les challengers} : comparer les modèles sur les mêmes prêts avant de
changer celui qui décide.
\end{enumerate}

\begin{piege}{Piège --- Écraser le score}
La solution naïve consiste à ajouter une colonne \code{score} dans la table des prêts et à la mettre à
jour à chaque nouveau modèle. On perd alors tout l'historique : impossible de savoir quel modèle a
accordé le prêt, impossible de comparer deux versions. Une prédiction est un \textbf{fait daté} : on
l'insère, on ne la met jamais à jour.
\end{piege}

Ta base conserve aussi le \textbf{découpage apprentissage / test} de ton projet R (table
\code{echantillons_modele}). Mesurer la performance d'un modèle sur les prêts qui ont servi à
l'entraîner donnerait des résultats trop beaux : on l'évaluera uniquement sur les 6\,482 prêts de test.

\subsection{Problème 4 --- Automatiser les alertes et le passage en défaut}

Avec 32\,409 prêts, personne ne peut vérifier chaque matin lesquels ont franchi 90 jours de retard.
Or un défaut détecté en retard fausse les provisions, les rapports au superviseur et les décisions de
recouvrement. Il faut donc que la base applique la règle \textbf{elle-même}.

Deux situations très différentes doivent être traitées :

\begin{itemize}
\item \textbf{Un événement se produit} : un paiement arrive. Le statut du prêt doit être recalculé
immédiatement : un client qui régularise redevient \og sain \fg{}. C'est le rôle d'un
\textbf{trigger}, une fonction que la base exécute automatiquement quand une ligne change.
\item \textbf{Rien ne se passe} : le client ne paie plus, et les jours défilent. Le retard franchit 90
jours \emph{sans qu'aucune ligne ne change}. Aucun trigger ne peut se déclencher : il faut un
\textbf{traitement planifié} (le « batch de nuit ») qui recalcule les statuts de tout le portefeuille.
\end{itemize}

\begin{regle}[Règle d'or de l'automatisation]
Un trigger réagit à un \textbf{événement} (une insertion, une modification, une suppression). Le
\textbf{passage du temps n'est pas un événement}. Toute règle qui dépend de la date du jour demande un
traitement planifié.
\end{regle}

Chaque changement de statut doit enfin laisser une trace : l'ancien statut, le nouveau, le retard à ce
moment, la date et l'utilisateur. C'est le \textbf{journal d'audit}. Et un passage en défaut doit
\textbf{alerter} le service de recouvrement : la base enverra une notification que toute application
abonnée recevra.

\begin{observation}[Pas de sortie automatique du défaut]
Un prêt en défaut qui régularise ne redevient pas sain immédiatement. Les orientations européennes sur
la définition du défaut imposent une \textbf{période probatoire} (au moins trois mois) avant d'en
sortir. Ton système appliquera donc une règle asymétrique : l'entrée en défaut est automatique, la
sortie demande une décision humaine.
\end{observation}

% -----------------------------------------------------------------------------
\section{Les objectifs techniques}

Les quatre problèmes métier se traduisent en quatre objectifs techniques, qui structurent la suite du
cours.

\begin{center}\small
\begin{tabularx}{\linewidth}{@{}B{3.3cm} L L P{2.2cm}@{}}
\toprule
Objectif & Concrètement & Notions SQL & Où \\
\midrule
Concevoir le schéma relationnel & 8 tables reliées, des données qui ne peuvent pas être incohérentes &
types, clés, \code{CHECK}, \code{UNIQUE}, index, transactions & Partie 3, Chapitre 1 \\
Ingérer des volumes & 1,3 million de lignes chargées en une quinzaine de secondes, en une seule
transaction & \code{COPY}, séquences, \code{ANALYZE} & Tâche 3 \\
Écrire du SQL analytique performant & retards, PAR, provisions, déciles, AUC, matrices de transition,
générations & vues, CTE, \code{FILTER}, fonctions de fenêtre, \code{ROLLUP} & Chapitre 2 \\
Automatiser la logique métier & recalcul des statuts, journal, alertes, encaissement, clôture mensuelle &
fonctions PL/pgSQL, triggers, procédures, rôles & Chapitre 3 \\
\bottomrule
\end{tabularx}
\end{center}

\subsection{L'architecture d'ensemble}

La figure~\ref{fig:architecture} montre comment les pièces s'assemblent. Les données sources et les
sorties de tes modèles R sont chargées par un script Python. La base contient les données \emph{et} la
logique métier. Plusieurs utilisateurs s'y connectent, chacun avec des droits différents.

\begin{figure}[ht]
\centering
\begin{tikzpicture}[font=\small, >=Stealth,
  boite/.style={draw=grisTitre, fill=white, rounded corners=3pt, align=center, minimum height=0.95cm, text width=3.3cm},
  conso/.style={draw=vertCadre, fill=vertFond, rounded corners=3pt, align=center, minimum height=0.9cm, text width=3.4cm}]
  \node[boite] (csv) at (0,1.4) {\textbf{Prêts réels}\\{\scriptsize credit\_risk\_dataset.csv}};
  \node[boite] (r) at (0,-0.2) {\textbf{Sorties des modèles R}\\{\footnotesize PD, SHAP, échantillons}};
  \node[boite, draw=orangeEx, fill=orangeFond] (py) at (0,-2.1) {\textbf{Chargeur Python}\\{\footnotesize Faker, Pandas, SQLAlchemy}};
  \draw[->, thick] (csv.west) -- ++(-0.35,0) |- (py.west);
  \draw[->, thick] (r) -- (py);
  % base
  \node[draw=bleuML, fill=bleuFond, rounded corners=4pt, line width=0.9pt, minimum width=5cm,
        minimum height=5.6cm] (db) at (5.3,-0.35) {};
  \node[font=\footnotesize\bfseries, color=bleuML] at ([yshift=-0.35cm]db.north) {PostgreSQL --- risk\_analytics\_db};
  \node[draw=bleuML!60, fill=white, rounded corners=2pt, text width=4.3cm, align=left, font=\footnotesize]
    (stock) at (5.3,0.75) {\textbf{Données}\\ 8 tables, contraintes, index\\ clients, prêts, échéances,\\
    prédictions, journal, provisions};
  \node[draw=bleuML!60, fill=white, rounded corners=2pt, text width=4.3cm, align=left, font=\footnotesize]
    (logique) at (5.3,-1.55) {\textbf{Logique métier}\\ 1 vue, 7 fonctions, 3 triggers,\\ 3 procédures,
    2 rôles};
  \draw[->, thick] (py.east) -- node[below, font=\footnotesize] {COPY} ([yshift=-1.75cm]db.west);
  % consommateurs
  \node[conso] (ana) at (10.9,1.75) {\textbf{Analyste risque}\\{\footnotesize DBeaver : lecture seule}};
  \node[conso] (app) at (10.9,0.45) {\textbf{Application paiement}\\{\footnotesize CALL enregistrer\_paiement}};
  \node[conso] (rep) at (10.9,-0.85) {\textbf{Comité, comptabilité}\\{\footnotesize provisions mensuelles}};
  \node[conso] (ds) at (10.9,-2.15) {\textbf{Data scientist}\\{\footnotesize monitoring des modèles}};
  \node[conso, draw=rouge, fill=rougeFond] (al) at (10.9,-3.4) {\textbf{Recouvrement}\\{\footnotesize alertes de défaut}};
  \foreach \c in {ana,app,rep,ds} \draw[<->, thick, grisTexte] (db.east |- \c) -- (\c.west);
  \draw[->, thick, rouge] (db.east |- al) -- node[above, font=\footnotesize] {NOTIFY} (al.west);
\end{tikzpicture}
\caption{L'architecture du système : des sources au pilotage du risque.}
\label{fig:architecture}
\end{figure}

\subsection{Les critères de réussite}

Comment sauras-tu que ton projet fonctionne ? Grâce à des \textbf{chiffres de contrôle}. Si tu
retrouves les valeurs ci-dessous, ton système est correct.

\begin{center}\small
\begin{tabularx}{\linewidth}{@{}L P{4.6cm} P{2.3cm}@{}}
\toprule
Contrôle & Valeur attendue & Où \\
\midrule
Prêts, mensualités et prédictions chargés & 32\,409 / 1\,107\,936 / 97\,227 & Tâche 3 \\
Prêts par statut après le premier traitement & 7\,088 en défaut, 933 en retard, 13\,480 sains,
10\,909 soldés & Chapitre 3 \\
Défauts détectés par ta logique \emph{contre} défauts réels & 7\,088 contre 7\,088, \textbf{0 désaccord} &
Chapitre 3 \\
PAR30 et PAR90 & 41,90~\% et 41,04~\% & Chapitre 2 \\
Provisions totales au 30/09/2026 & environ 25,63 millions d'euros (19,13~\% de l'EAD) & Chapitres 2 et 3 \\
AUC et Gini du champion sur l'échantillon de test & 0,923 et 0,847 & Chapitre 2 \\
Retard d'un prêt, calculé à la demande & en moins d'une milliseconde & Chapitre 1 \\
\bottomrule
\end{tabularx}
\end{center}

\begin{resume}
\begin{enumerate}
\item Le \textbf{DPD} se mesure depuis l'échéance impayée \textbf{la plus ancienne} ; il se calcule, il
ne se stocke pas. Les tranches de retard donnent le PAR30 et le PAR90, en montant.
\item L'\textbf{EAD} est le capital restant dû. La provision vaut $PD \times LGD \times EAD$, avec une PD à
12 mois (stage 1), sur la durée restante (stage 2) ou égale à 100~\% (stage 3) : $PD_m = 1-(1-PD)^{m/n}$.
\item Les \textbf{prédictions} sont des faits datés : on les insère dans un registre, on ne les modifie
jamais. Le registre sert à l'audit, à l'explication, au monitoring et à la comparaison des modèles.
\item Un \textbf{trigger} réagit aux événements, un \textbf{traitement planifié} gère le passage du
temps. L'entrée en défaut est automatique, la sortie est humaine.
\item Quatre objectifs techniques : schéma, ingestion, SQL analytique, automatisation.
\end{enumerate}
\end{resume}
